%%%%%%%%%%%%%%%%%%%%%%%%%%%%%%%%%%%%%%%%%
% Wnioski do pracy dyplomowej
% Szablon pracy dyplomowej
% Wydział Informatyki 
% Zachodniopomorski Uniwersytet Technologiczny w Szczecinie
% autor Joanna Kołodziejczyk (jkolodziejczyk@zut.edu.pl)
% Bardzo wczesnym pierwowzorem szablonu był
% The Legrand Orange Book
% Version 5.0 (29/05/2025)
%
% Modifications to LOB assigned by %JK
%%%%%%%%%%%%%%%%%%%%%%%%%%%%%%%%%%%%%%%%%


\chapter*{Podsumowanie}

% -------------------------------------------------------
% 1. Przypomnienie celu
% -------------------------------------------------------
Celem niniejszej pracy było zaprojektowanie i implementacja interfejsu
mózg-komputer umożliwiającego sterowanie kursorem myszy za pomocą
wzrokowych potencjałów wywołanych stanu ustalonego (SSVEP).
Postawione wymagania obejmowały obsługę pięciu komend sterujących
(ruch w czterech kierunkach oraz kliknięcie lewym przyciskiem myszy),
działanie w czasie rzeczywistym oraz łączne opóźnienie systemu
nieprzekraczające 5 sekund.

% -------------------------------------------------------
% 2. Co wykonano
% -------------------------------------------------------
W ramach pracy przeprowadzono analizę literatury dotyczącej interfejsów
BCI opartych na SSVEP, ze szczególnym uwzględnieniem metod klasyfikacji
sygnału EEG. Na tej podstawie zaprojektowano i zaimplementowano system
w języku Python, składający się z modułu akwizycji sygnału EEG
(płytka OpenBCI Cyton, biblioteka BrainFlow), modułu przetwarzania
sygnału (filtracja FIR z oknem Hamminga, filtr zaporowy 50~Hz,
biblioteka MNE-Python) oraz modułu klasyfikacji opartego na analizie
korelacji kanonicznej (CCA, biblioteka scikit-learn). Moduł stymulujący
zrealizowano przy użyciu biblioteki PyQt5 w postaci pełnoekranowej
nakładki z pięcioma migającymi bodźcami o częstotliwościach 6,66~Hz,
7,5~Hz, 8,57~Hz, 10,0~Hz i 12,0~Hz, dobranych jako całkowite
dzielniki częstotliwości odświeżania monitora (60~Hz).
Architektura wieloprocesowa z kolejką komend zapewnia niezależność
modułu stymulacji od obliczeń klasyfikatora.
Przeprowadzono również badanie z udziałem trzech uczestników,
obejmujące zadanie sterowania kursorem oraz zadanie skuteczności
klasyfikacji komend, uzupełnione ankietą subiektywnej oceny systemu.

% -------------------------------------------------------
% 3. Główne wyniki
% -------------------------------------------------------
Zrealizowany system spełnia postawione wymagania funkcjonalne
i pozafunkcjonalne. Minimalne opóźnienie od momentu skupienia wzroku
do wykonania komendy wynosi 4,2~s, mieszcząc się w założonym limicie
5~s. Badanie z udziałem trzech uczestników wykazało zróżnicowane,
lecz obiecujące wyniki. Dwóch uczestników osiągnęło skuteczność
klasyfikacji komend na poziomie odpowiednio 87\% i 100\%,
przy czym u drugiego uczestnika zaobserwowano problem wielokrotnych
aktywacji tej samej komendy, wynikający z silnej odpowiedzi SSVEP
przekraczającej próg klasyfikatora przez cały czas skupienia wzroku.
Trzeci uczestnik wykazał znaczące trudności z obsługą systemu
przy jednoczesnej obecności pięciu bodźców na ekranie,
osiągając skuteczność 60\% przy 33\% błędnych klasyfikacji.
Zjawisko to jest zgodne z opisywanym w literaturze fenomenem
\textit{BCI inefficiency}, dotyczącym około 15--25\% populacji
\cite{vialatte2010}. Czasy ukończenia zadania sterowania kursorem
wyniosły średnio 281~s i 236~s dla dwóch uczestników,
którzy zdołali je ukończyć.

% -------------------------------------------------------
% 4. Wkład w dziedzinę
% -------------------------------------------------------
Praca potwierdza praktyczną przydatność metody CCA jako klasyfikatora
niewymagającego kalibracji per użytkownik w systemach SSVEP
działających w czasie rzeczywistym. Wyniki badania wskazują,
że nawet przy relatywnie długim oknie analizy (4~s) system
jest zdolny do stabilnej klasyfikacji u większości użytkowników.
Zidentyfikowano dwa istotne ograniczenia implementacyjne:
brak adaptacyjnego progu klasyfikacji prowadzący do wielokrotnych
aktywacji u użytkowników z silną odpowiedzią SSVEP
oraz brak mechanizmu selekcji częstotliwości stymulacji
dopasowanego do indywidualnych cech użytkownika.

% -------------------------------------------------------
% 5. Kierunki dalszych badań
% -------------------------------------------------------
Dalszy rozwój systemu mógłby objąć wprowadzenie krótkiej sesji
kalibracyjnej wyznaczającej indywidualny próg $\theta$ oraz
optymalny czas opóźnienia wykonania kolejnych komend dla każdego użytkownika.
Wartościowym uzupełnieniem byłoby również dodanie wizualnego
wskaźnika bieżącej wartości współczynnika korelacji kanonicznej
$\rho$ w interfejsie graficznym, co pozwoliłoby użytkownikowi
na bieżąco oceniać jakość skupienia wzroku i skrócić czas reakcji.
Na podstawie sugestii uczestników badania warto rozważyć eksperymenty z bodźcami wyświetlanymi na tle pulpitu systemu operacyjnego zamiast jednolitego czarnego ekranu. Takie podejście mogłoby zmniejszyć dyskomfort wzrokowy przy dłuższym użytkowaniu, choć wiązałoby się z ryzykiem obniżenia amplitudy odpowiedzi SSVEP
Rozszerzenie zestawu elektrod o dodatkowe kanały potyliczne
(POz, PO3, PO4) mogłoby zwiększyć czułość klasyfikatora
u uczestników z słabą odpowiedzią SSVEP.

Interesującym kierunkiem byłoby także zbadanie możliwości
zastąpienia lub uzupełnienia metody CCA algorytmem
FBCCA (Filter Bank CCA), który poprzez analizę sygnału
w wielu podzakresach częstotliwościowych wykazuje wyższą
skuteczność detekcji przy krótszych oknach analizy \cite{chen2015}.